\centerline{\Huge \bf  Тренировочная Олимпиада 5}
\centerline{\Huge \bf 8 класс}
\centerline{\Huge \bf Уровень: Заключительный отбор в Сириус}

\begin{flushright}
    \scriptsize{Глава 5. Вы двое, танцуйте, если хотите победить!}
\end{flushright}

\problem{1}
На доске выписаны в ряд 5 целых чисел. Эрдни утверждает, что сумма любых двух подряд идущих отрицательных, а Соня утверждает, что суммы любых трёх подряд идущих положительна. Докажите, что кто-то из них ошибается.
\begin{flushright}
    \textit{\small{(Гасанов Р.В.)}}
\end{flushright}

\problem{2}
РВ и АА играли в морской бой, где у каждого имеется поле $14 \times 29$ и по одному кораблю размером $2\times 2$. АА решил считерить и первым же ходом нанёс 49 баллистических ударов, каждый из которых поражает область $2\times 2$. Есть ли шансы на выживание у корабля РВ?

\textit{Примечание:} Корабль считается выжившим, если в него вообще не задели.
\begin{flushright}
    \textit{\small{(Гасанов Р.В.)}}
\end{flushright}

\problem{3}
РВ и ОВ решили составить олимпиады на много лет вперед. Они записывают в свои листки задачи, и когда на листке появляется 5 задач, то лист становится олимпиадным, и как отдельная олимпиада добавляется в архив. РВ составляет задачи в 1.5 раза быстрее ОВ. В какой-то не начальный момент времени на листке перед ОВ и листке перед РВ было одинаковое количество задач. РВ понял, что больше так не может и пошел спать. ОВ же продолжил составлять олимпиады. Через 4 часа, когда РВ выспался, он заметил, что ОВ обогнал его ровно на одну олимпиаду. Сколько задач было на листке у РВ перед тем как он пошел спать, если ОВ составлял задачи с минимально возможной скоростью?
\begin{flushright}
    \textit{\small{(Очиров О.В.)}}
\end{flushright}

\problem{4}
Найдите геометрическое место точек, сумма квадратов расстояний от которых до вершин острых углов равнобедренного прямоугольного треугольника вдвое больше квадрата расстояния до вершины прямого угла.
\begin{flushright}
    \textit{\small{(Гасанов Р.В.)}}
\end{flushright}


\newpage

\hspace{3mm}

\problem{5}
РВ и Герман играют в дагестанскую рулетку. Правила следующие: есть один шестизарядный револьвер с одним патроном. Сначала РВ крутит барабан и стреляет в воздух. Если выстрела не произошло, то ход переходит Герману и он делает то же самое. Тот у кого револьвер выстрелил бросает на прогиб своего опонента. Какова вероятность, что Герман будет брошен на прогиб?

\begin{flushright}
    \textit{\small{(Гасанов Р.В.)}}
\end{flushright}

\hspace{3mm}

\problem{6}
Имеет ли решение уравнение 
\[ \dfrac{1}{n} + \dfrac{1}{n+1} + \ldots + \dfrac{1}{n^2 - 1} + \dfrac{1}{n^2} = 1\]
в натуральных числах, кроме $n = 1$?
\begin{flushright}
    \textit{\small{(Гасанов Р.В.)}}
\end{flushright}
